% ============================================================
%  chapters/05_conclusion.tex
% ============================================================

\section{Conclusion}
\label{sec:conclusion}

Evidence from 78~histopathology AI fairness studies (2017--2026) reveals a sharp and uncomfortable truth: fairness in computational pathology is empirically measurable and methodologically tractable, yet remains structurally unmeasured. The field has built the instruments to detect its own failures. TSS leakage is recoverable at $>$0.96~AUROC across multiple TCGA cancer types and at greater than~90\%~linear-probe accuracy from state-of-the-art foundation-model embeddings. Demographic subgroup AUROC gaps of 3--16\% persist in breast, lung, and colorectal tasks, with 10--25\% cross-population performance degradation at geography boundaries.

Individual fairness interventions have already demonstrated clinically meaningful effects. Knowledge-guided fine-tuning achieves 61\% demographic-gap reduction, fairness-aware contrastive objectives deliver 88.5--91.1\% disparity mitigation across 20 TCGA cancer types (88.5\% internal, 91.1\% external validation; range reflects validation setting), conditional diffusion yields 48.5\% improvement in demographic parity, and SPARE improves minority-group accuracy by 3.7--5.8\%. Proportional federated averaging reduces inter-hospital performance variance from~32.54 to~10.00, and causal survival models equalise recurrence-free-survival performance across seven geographically distributed sites. These gains remain at scales insufficient to bound aggregate TSS deployment risk.

Yet only 19.2\%~of studies report any fairness metric, 7.7\% report worst-group performance, and 29.5\% (95\% CI: 20.5--40.2\%) perform external validation. The geographic concentration is pronounced: 84.6\% (66/78) of first/corresponding authors are from North America and Europe, with East Asia contributing most of the remainder. Meanwhile, review recommendations consistently endorse demographic reporting ($\sim$90\%), multi-site evaluation ($\sim$85\%), and preserved-site cross-validation ($\sim$90\%), yet experimental practice fails to supply these. This 50--75~percentage-point gap between review consensus and experimental practice reflects not an epistemic failure but an infrastructural one. The checklist, benchmark, methodological priorities, and policy levers proposed in this review specify precisely the reporting, evaluation, representational, and incentive substrates whose absence the preceding sections have quantified. The state of the field in 2026 is therefore an enforcement problem, not a discovery problem.

Closing this gap requires coordinated action across four constituencies. Researchers must adopt preserved-site cross-validation, subgroup-stratified reporting, and embedding-level TSS leakage probes as minimum-viable evidence, and must publish datasets, code, and evaluation harnesses with per-field demographic completeness disclosed rather than implied. Clinicians, pathologists, and laboratory directors must treat scanner calibration, stain-protocol drift, and subgroup-stratified quality control as first-class workflow instruments and demand failure-mode transparency from deployed devices. Regulators and editors must bind the reporting checklist proposed in this review to pre-market authorisation and journal-submission policy, and require post-market subgroup surveillance with the same rigour already imposed on aggregate performance. Funders, consortia, and patient communities must resource the multi-continental, demographically stratified cohorts and the federated, consented governance infrastructure without which fairness-aware pretraining, causal estimation, and synthetic augmentation remain exercises on a distribution that no longer adequately represents the populations computational pathology will be asked to serve.

Each required artefact (the checklist, the benchmark, the mitigation programme, and the incentive layer) lies within reach of the community that has already produced the evidence documented in this review; none will emerge from exhortation alone. The next decade of histopathology AI will be judged not by whether its models achieve higher accuracy on aggregate benchmarks, but by whether the field chooses to instrument, report, and govern the subgroup behavior of the systems it is now placing into clinical practice. Equitable computational pathology is, on the evidence assembled here, an operational commitment rather than an aspirational one. The window to make that commitment on the field's own terms, rather than retrofitting it under regulatory, litigation, or public-trust pressure, is open now and will not remain open indefinitely.
